% Implementation, Oversight, Enforcement, and Fiscal Provisions. Built by the
% future pass from the instructions below. Do not process now. Covers the
% implementation choices, the oversight and enforcement authority, the phased
% rollout, and the budget, funding, and offsets. Carry the dense numbers in
% tables sourced from the national-platform financial and implementation sections.

[Write as connected legislative prose with two tables (an implementation-choices
table and a fiscal table). Synthesise the sources; do not narrate them. Movement 1
sets out the implementation design choices, Movement 2 the oversight and
enforcement authority, Movement 3 the phased rollout and workforce, and Movement 4
the budget, funding, and offsets.]

[Movement 1, Implementation design choices. State, with one mapping table, the
implementation levers the bill selects from the standard menu: mandatory
disclosure of AI use in robot-patient interactions; auditability through a hash-
chained audit trail and the five-server MCP-PAI topology; retained human review of
medical-necessity decisions; an FDA review touchpoint consistent with the device
and credibility framework; a CMS reimbursement pathway for verified Physical AI
procedures; a grant program seeding site readiness; and a periodic reporting duty.
Source the infrastructure and governance choices from the national-platform
\texttt{national\_mcp.tex} (the five servers, 23 tools, role-based access, 21 CFR
Part 11 alignment) and \texttt{implementation\_strategy.tex} (standards-integration
gates, quality and risk management). Cite \texttt{kawchak2026mcp},
\texttt{cfr-part11}, \texttt{fda2023credibility}, and \texttt{usc-1395y}.]

[Movement 2, Oversight and enforcement authority. Designate the oversight roles
(an FDA touchpoint for device and credibility review; a state insurance or health
authority for the utilization-review-adjacent duties) and the enforcement
mechanics: audits and compliance reviews, civil penalties for willful violations
on the California model (\texttt{ca-sb1120}), a commissioner finding of
noncompliance triggering algorithm production on the Texas model
(\texttt{tx-sb1822}), and certification of bias minimization on the New York model
(\texttt{ny-a9149}). Make the operative trigger the absence of a cleared
verification-before-generation record. Cite \texttt{fl-hb527} for the human-review
backstop.]

[Movement 3, Phased rollout and workforce. Summarise a three-phase national
rollout (a first state, then multi-state, then nationwide) and the workforce roles
(Physical AI operators, technicians, safety specialists) from the national-
platform \texttt{implementation\_strategy.tex} timeline and workforce tables;
reference those tables rather than reproducing every row. Tie site activation to
the PSL gate and robot admission to the USL minimum. Cite \texttt{kawchak2026national}
and \texttt{kawchak2026sitedocs}.]

[Movement 4, Budget, funding, and offsets. Present the fiscal case in one table:
the cost of trial delay, the per-patient cost comparison, the single-site and
national projections, and the patient and sponsor benefits, drawn from the
national-platform \texttt{financial\_analysis.tex}. State the central offset
argument, that automating the VVUQ assurance layer is inexpensive relative to its
decision weight, supported by the VVUQ-01 and VVUQ-02 cost evidence (an autonomous
agent carried the heavier assurance half cheaply and reproducibly) from
\texttt{VVUQ-02/execution/README.md} (cost section) and the VVUQ-01 cost-and-
credibility bridge in \texttt{VVUQ-01/final-paper/sections/discussion.tex}. Frame
the savings from faster, safer trials as the funding offset for the grant program
and oversight costs. Keep all figures as planning indices or simulation results,
not guarantees. Cite \texttt{kawchak2026vvuq01}, \texttt{kawchak2026vvuq02}, and
\texttt{fda2026realtime}.]

% Model implementation-choices table (parent paths at top, abbreviated leaves):
% \begin{table}[h]\centering
% \begin{tabularx}{\textwidth}{L{3.4cm} Y}
% \toprule
% \textbf{Lever} & \textbf{Bill choice and source} \\
% \midrule
% \multicolumn{2}{l}{\textit{Parent: physical-ai .../final\_paper/sections/}} \\
% Auditability & Hash-chained trail, five MCP servers (national\_mcp.tex) \\
% Human review & Retained medical-necessity decision (regulatory\_landscape.tex) \\
% Reimbursement & CMS pathway for verified procedures (financial\_analysis.tex) \\
% Rollout & Three phases, workforce roles (implementation\_strategy.tex) \\
% \bottomrule
% \end{tabularx}
% \caption{Implementation levers selected by the bill and their sources.}
% \end{table}
